\begin{abstract}
传统辐射源定位通常先估计站间到达时间差，再求解源位置。当多个未知波形占用同一频带时，混合信号增加了时延提取和跨站关联的难度。直接位置估计（DPD）利用多站观测，在统一模型中联合估计源位置。然而，强弱源共存时，粗搜索中排名靠前的位置未必对应精修后的最佳多源解；逐源精修若固定信号系数，又可能因位置与波形估计的耦合而产生大量小步迭代。针对这两个求解问题，本文研究同步单天线站在已知自由空间相对幅度与基带时延模型下的未知宽带双源定位。首先以最小二乘消去未知波形，得到仅关于位置的优化目标；在此基础上，提出多候选联合精修交替投影（AP）与波形消元块NOMP。前者保留条件谱中的少量候选，分别联合调整两个源的位置，按最终拟合残差择优；后者在残差发现与循环精修中，对每个试探位置重新估计信号系数，再计算位置更新方向。条件投影与曲率推导说明了两种策略的作用。已完成的120条双源配对观测中，多候选联合AP恢复了单候选遗漏的10条观测；相同停止精度下，单候选联合精修的平均求解时间比交替精修减少约8.2\%。八条观测的受控实验进一步表明，波形消元减少了固定系数更新中的重复小步计算。结果支持从候选选择与局部更新两个环节改善强弱重叠源定位的可靠性与计算效率。
\end{abstract}
\noindent\textbf{关键词：}直接位置估计；强弱重叠源；交替投影；块NOMP；可分离非线性最小二乘
